\section{Relation to existing work}
Electric-bus scheduling combines service assignment, battery feasibility and
charging access. Grid-aware multi-depot scheduling, capacitated stations and
scalable route-and-recharge generation are studied by
\citet{wu2022grid,devos2024capacity,parmentier2023fleets}. These operational
problems motivate our use of a complete fleet: charging feasibility is joint
across the selected vehicle duties, especially at a shared connector.
Our question concerns support over alternatives that can change both the
number of used buses and the charging opportunities created by their duties.
Integrated flow and set-partitioning models already combine assignments,
battery feasibility and charging slots \citep{borndorfer2024integrated}.
\citet{wu2022grid} already choose vehicle-to-trip assignments and charging
under grid-related objectives; full overnight charging is among their
assumptions. Endogenous assignments, full charging and grid-aware scheduling
are therefore not, separately or together, sufficient novelty claims here.

\citet{zoltowskaLin2021} formulate aggregation and disaggregation of bus
charging around hourly auction bids using integer models. Their work is a
precedent for discrete bus charging and market participation; our question
additionally allows service assignments and used-bus costs to change under
explicit terminal replacement-energy accounting. More directly,
\citet{wu2019gridinteractive} describe bus service planning coupled to market
clearing and distribution locational marginal prices in a bilevel framework.
Our question is the support of complete physical fleets at a fixed posted
marginal tariff, rather than a first coupling of bus operations to endogenous
electricity prices.

Price-support failure with indivisible operating decisions is established
in electricity-market research. \citet{oneill2005} study prices in markets
with nonconvexities; \citet{gribikHoganPope2007} discuss market-clearing
prices and uplift. Convex-hull pricing provides the relevant convexification
framework \citep{schiro2016,huaBaldick2017}. Linear-price support in
decentralized problems with integralities is also studied through complete
linear descriptions \citep{huembs2022complete}. We use this machinery; we do
not present the support equivalence below as a new general theorem.
Whole feasible participant schedules as columns and their fixed-price
self-scheduling oracles also have established decomposition treatments
\citep{andrianesis2022hull}. The complete-fleet unit ensures the right
physical alternatives in this application; it is not new column-generation
machinery.

Exact EV flexibility aggregation is also established.
\citet{mukhiAbate2023exact} characterize fixed-population charging polytopes
using permutahedra, with prescribed availability and full departure charge.
\citet{mukhiLohoAbate2025exact} give generalized-polymatroid representations
and exact Minkowski aggregation for convex individual flexibility sets.
Our distinction is the joint load--vehicle-cost frontier when service pairing
and the number of used buses change. The load-only projection in our class
is an interval, not a new nonconvex flexibility set. Fixed-matching interval
aggregation and ordered matching are classical ingredients.

Bidirectional charging and transport decisions have also been combined
previously. \citet[Sections 2.1--2.2]{blatiak2022v2g} optimize commercial
fleet trip timing and charging against energy and balancing-service prices.
The public author manuscript of \citet[Section 2]{zhou2025v2g} jointly
chooses bus use, trip connections and partial charging/discharging under
input purchase and export tariffs. Its objective includes vehicle fees,
deadheading, charging, export revenue and degradation, with full-depot
inventory and overnight recharging described in the model and case study.
These features materially overlap our physical setting. Our comparison
uses the author manuscript; imperfectly extracted mathematical glyphs do
not justify detailed claims about its connector or terminal constraints.
The distinction is our exact cost-labelled threshold and own-load support
classification, not a richer operational V2G formulation.
Neither coupling transport to V2G nor using spare buses for grid services
is a novelty claim here.

Signed aggregate flexibility is established as well.
\citet[Section II, Example 1]{altaha2024flexibility} model lossless
bidirectional batteries with initial and final energy requirements and
aggregate a fixed population through Minkowski sums. Our V2G extension
instead classifies the joint load--fleet-cost frontier created by successor
matching and a shared gross-throughput budget. Its thresholds and paid-idle
elimination argument concern that restricted physical class; the
persistence and restoration examples illustrate the classification, rather
than establishing a first signed-aggregation or V2G price-support result.

Lossy storage also has exact coordinate reformulations.
\citet[Theorems 1--3]{altahaBitar2024lossy} establish a piecewise-affine
grid-power/SOC bijection, a convex SOC feasible set for a single store,
and sufficient monotonicity conditions for a convex reformulated objective.
Our common-loss result does not claim these transformations or convexity
conditions as new. It combines service-dependent vehicle counts and
successor choices with shared gross connector use to describe the efficient
joint early-load, terminal-load and vehicle-cost envelope. Its finite
piecewise description is conditional on the stated common-window model,
and excludes dominated points only under explicitly monotone economics.

Aggregation must specify who owns the aggregated decisions. The classical
perspective of \citet{starr1969} and dispatchable-demand results of
\citet{hreinsson2021} explain why aggregate nonconvexities can become small.
Our replication calculation instead isolates a distinction between a
planning-value gap, one owner's absolute fixed-price incentive and the
incentives of independently reserved service pairs. The explicit rates do
not contradict approximate aggregation results.

The comparison concerns what can be decided from an explicit frontier,
rather than whether assignments or storage can be optimized at all:
\begin{center}
\small
\begin{tabular}{p{0.40\textwidth}p{0.52\textwidth}}
\toprule
Established viewpoint & Decision supplied by this classification\\
\midrule
Operational fleet models choose assignments and charging for a specified
objective \citep{wu2022grid,borndorfer2024integrated,zhou2025v2g}.
& Theorem~\ref{thm:sharedprojection} gives every attainable load--vehicle-cost
layer in the stated two-batch class; Corollary~\ref{cor:sharedsupport}
tests physical attainment of its hull optimum.\\
Fixed-population flexibility represents aggregate battery loads
\citep{mukhiAbate2023exact,altaha2024flexibility}.
& Changing successor matches and vehicle counts changes operating-cost
layers even when the load-only set is an interval. The thresholds retain
this distinction.\\
Lossy-storage reformulation relates grid power and inventory
\citep{altahaBitar2024lossy}.
& Theorem~\ref{thm:lossyfrontier} reduces the assignment-dependent efficient
joint load--cost set to at most two segments per cardinality, preserving
planning values under monotone supply costs.\\
\bottomrule
\end{tabular}
\end{center}
These are restricted structural decisions, not claims of a richer operational
model or faster solution of general bus scheduling. Standard matching
integrality, value-function convexity and price-support duality remain the
underlying tools. The capacity, fee and equipment examples explain the
classification's consequences; the separately certified public instance
establishes existence within its declared movement graph, rather than
extending the closed-form theorem to arbitrary timetables.
